\section{\label{sec:spdf}与光前分布的关系}

我们通过考察这两类分布的 Mellin 矩，并利用 Mellin 矩与局域算符矩阵元之间的联系——对光前 PDFs 而言这些算符是扭度2的——来讨论准 PDF 与光前 PDF 的关系 \cite{Christ:1972ms}。对准 PDFs 而言，与 Mellin 矩对应的局域算符没有良定义的扭度，但在扣除高扭度效应并施加靶质量修正之后，它们可与扭度2算符联系起来
\cite{Lin:2014zya,Alexandrou:2015rja}。虽然本文考虑的是涂抹矩阵元，文献 \cite{Lin:2014zya,Alexandrou:2015rja} 中关于高扭度与靶质量效应的论证仍然适用，
因为流时间充当着替代格距的规范不变调节子。

我们按如下方式把准 PDF 的 Mellin 矩与局域算符的矩阵元联系起来。在轴规范 $B_z(x;\tau)=0$ 下，
矩阵元 $h^{(s)}$ 为
\begin{equation}\label{eq:hAzdef}
h^{(s)}\left(\frac{z}{\sqrt{\tau}},\sqrt{\tau}P_z,\sqrt{\tau}\Lambda_{\mathrm{QCD}},
\sqrt{\tau}
M_\mathrm{N}\right)\Big|_{B_z=0} =
\frac{1}{2P_z}\left\langle
P_z \left| \overline{\chi}(z;\tau)
\gamma_z\frac{\lambda^a}{2}\chi(0;\tau)\right|
P_z\right\rangle_\mathrm{C}.
\end{equation}
现将该表达式代入准 PDF 的定义，即式~\eqref{eq:qpdfdef}，并把所得表达式在 $\xi$ 的整个范围上积分。与光前
PDF 不同，该范围从负无穷延伸到正无穷，于是得到
\begin{equation}
\int_{-\infty}^\infty
\mathrm{d}\xi\,
q^{\,(s)}\left(\xi,\sqrt{\tau}P_z,\sqrt{\tau}\Lambda_{\mathrm{QCD}},\sqrt{\tau}
M_\mathrm{N}\right)\Big|_{B_z=0}
=
h^{(s)}(0,\sqrt{\tau}P_z,\sqrt{\tau}\Lambda_{\mathrm{QCD}},\sqrt{\tau}M_\mathrm{
N}
)\Big|_{B_z=0}.
\end{equation}
这里利用了关系 $\delta(zP_z) =
\delta(z)/P_z$（对 $P_z>0$）。我们看到，准 PDF 的第一个 Mellin 矩可以用一个局域
（涂抹）算符的欧几里得矩阵元表出。

通过考虑准分布对空间间隔 $z$ 的导数，可以把这个论证推广到准 PDFs 的任意阶矩 \cite{Collins:2011zzd}。对式~\eqref{eq:qpdfdef} 作傅里叶逆变换，得
\begin{equation}\label{eq:handqpdf}
h^{(s)}\left(\frac{z}{\sqrt{\tau}},\sqrt{\tau}P_z,\sqrt{\tau}\Lambda_{\mathrm{QCD}},
\sqrt{\tau}
M_\mathrm{N}\right)
=
\int_{-\infty}^\infty \mathrm{d}\xi\, e^{-i\xi z
P_z}
q^{\,(s)}\left(\xi,\sqrt{\tau}P_z,\sqrt{\tau}\Lambda_{\mathrm{QCD}},\sqrt{\tau}
M_\mathrm{N}\right).
\end{equation}
再对位移 $z$ 求导数，得
\begin{align}\label{eq:hint1}
\left(\frac{i}{P_z}\frac{\partial}{\partial
z}\right)^{n-1}
h^{(s)}{} &
\left(\frac{z}{\sqrt{\tau}},\sqrt{\tau}P_z,\sqrt{\tau}\Lambda_{\mathrm{QCD}},
\sqrt{\tau}
M_\mathrm{N}\right) = \nonumber\\
{} & \int_{-\infty}^\infty \mathrm{d}\xi\, \xi^{n-1}e^{-i\xi z
P_z}
q^{\,(s)}\left(\xi,\sqrt{\tau}P_z,\sqrt{\tau}\Lambda_{\mathrm{QCD}},\sqrt{\tau}
M_\mathrm{N}\right).
\end{align}

把涂抹准 PDF 的矩 $b_n^{(s)}$ 定义为
\begin{equation}\label{eq:bdef}
b_n^{(s)}\left(\sqrt{\tau}P_z,\frac{\Lambda_{\mathrm{QCD}}}{P_z},
\frac{M_\mathrm{N}}{P_z}\right)
= \int_{-\infty}^\infty
\mathrm{d}\xi\, \xi^{n-1}
q^{\,(s)}\left(\xi,\sqrt{\tau}P_z,\sqrt{\tau}\Lambda_{\mathrm{QCD}},\sqrt{\tau}
M_\mathrm{N}\right),
\end{equation}
并把式~\eqref{eq:hdef} 给出的矩阵元 $h^{(s)}$ 的定义代入式~\eqref{eq:hint1}，在
$z\rightarrow 0$ 的极限下便得到
\begin{align}
b_n^{(s)}\left( \sqrt{\tau}P_z,\frac{\Lambda_{\mathrm{QCD}}}{P_z},
\frac{M_\mathrm{N}}{P_z}\right)_{B_z=0}
= {} & \frac{c_n^{(s)}(\sqrt{\tau}P_z)}{2P_z^n} \nonumber\\
{} & \qquad \times \left\langle
P_z \left| \left[\overline{\chi}(z;\tau) \gamma_z
\left(i\overleftarrow{\partial}_z^{n-1}\right)\frac{\lambda^a}{2}
\chi(0;\tau)\right ] _{z=0 } \right|P_z\right\rangle_\mathrm{C}.
\end{align}
微扰系数 $c_n^{(s)}(\sqrt{\tau}P_z)$ 刻画了式~\eqref{eq:hint1} 右端在间隔 z 与流时间 $\tau$ 同时趋于零时的潜在奇异性；如文献
\cite{Lin:2014zya} 所概述，它们可由处理非局域矩阵元的涂抹算符乘积展开
\cite{Monahan:2015lha} 得到。

恢复规范不变性后，我们得到准 PDFs 之矩用局域算符欧几里得矩阵元表示的最终表达式：
\begin{equation}
b_n^{(s)} \left(\sqrt{\tau}P_z,\frac{\Lambda_{\mathrm{QCD}}}{P_z},
\frac{M_\mathrm{N}}{P_z}\right)
=\frac{c_n^{(s)}(\sqrt{\tau}P_z)}{2P_z^n}\left\langle
P_z \left| \left[
\overline{\chi}(z;\tau)\gamma_z (i\overleftarrow{D}_z)^{(n-1)}
\frac{\lambda^a}{2}\chi(0;\tau) \right ] _{z=0 }
\right|P_z\right\rangle_\mathrm{C} .
\end{equation}
该表达式右端矩阵元中出现的局域算符并非扭度2算符。准 PDF 之矩与扭度2算符矩阵元之间的差别由按
$\Lambda_\mathrm{QCD}^2/{P_z^2}$
与 $M_\mathrm{N}^2/P_z^2$ 标度的项给出 \cite{Lin:2014zya,Alexandrou:2015rja}。按 ${\cal O}(M_\mathrm{N}^2/P_z^2)$ 标度的项对应靶质量
修正 \cite{Nachtmann:1973ll,Georgi:1974sr}。尽管存在这些靶质量修正时 PDFs 的恰当诠释颇为微妙
\cite{Steffens:2012jx,Monfared:2014nta}，就本文目的而言只需知道，这些
非领头
修正可以通过改写而被吸收
\cite{Lin:2014zya,Alexandrou:2015rja}
\begin{align}
b_n^{(s)} \left(\sqrt{\tau}P_z,\frac{\Lambda_{\mathrm{QCD}}}{P_z}\right)
= {} & \frac{c_n^{(s)}(\sqrt{\tau}P_z)}{2P_z^n}\left\langle
P_z \left|\left[
\overline{\chi}(z;\tau)\gamma_z (i\overleftarrow{D}_z)^{(n-1)}
\frac{\lambda^a}{2}\chi(0;\tau) \right ] _{z=0 }
\right|P_z\right\rangle_\mathrm{C} \nonumber \\
{} & \qquad \times K_n^{-1}\left(\frac{M_\mathrm{N}^2}{4P_z^2}\right),
\end{align}
其中
\begin{equation}
K_n\left(\frac{M_\mathrm{N}^2}{4P_z^2}\right) = \sum_{j=0}^{n/2}
\bigg(
\begin{array}{c}
n-j \\
j
\end{array}
\bigg)
\left(\frac{M_\mathrm{N}^2}{4P_z^2}\right)^j.
\end{equation}

经此修正后的矩阵元现在可以对
$\Lambda_\mathrm{QCD}^2/P_z^2$ 作泰勒级数展开。该展开中的系数
代表高扭度效应，其起因在于原矩阵元并非扭度2算符的矩阵元。展开中的
领头
系数是扭度2贡献，只能依赖
于核子结构与流时间：
\begin{equation}\label{eq:btwist}
b_n^{(s)} \left(\sqrt{\tau}P_z,\sqrt{\tau}\Lambda_{\mathrm{QCD}}\right)
=
c^{(s)}_n(\sqrt{\tau}P_z)b_n^{(s,\mathrm{twist}-2)}
\left(\sqrt{\tau}\Lambda_{\mathrm{QCD}}\right)  +{\cal
O}\left(\frac{\Lambda_\mathrm{QCD}^2}{P_z^2}\right),
\end{equation}
其中当 $\Lambda_\mathrm{QCD}^2/P_z^2\ll 1$ 时，高扭度修正
可以忽略。

总结起来，我们假定：第一，可以对靶质量
修正作精确改正；第二，可以把动量 $P_z$ 取得足够大以致高
扭度效应可忽略。于是，在这些
假设下，涂抹准 PDFs 的矩成为
微扰系数与纯扭度2矩阵元的无量纲乘积，后者只是无量纲量
$\sqrt{\tau} \Lambda_{\mathrm{QCD}}$ 的函数，其中含有
关于强子结构的
信息。

\subsection{\label{sec:renormpdf} 短距离展开}

现在，借助短距离展开赋予梯度流的性质
\cite{Luscher:2011bx,Suzuki:2013gza,Asakawa:2013laa,Makino:2014taa,
Hieda:2016lly}，我们可以把涂抹准 PDF 的矩
$b_n^{(s,\mathrm{twist}-2)}\big(\sqrt{\tau}\Lambda_{\mathrm{QCD}}\big)$——它们是涂抹场的局域矩阵元——与
光前 PDFs 的重正化矩联系起来。
涂抹程序的指数定域性质
使得涂抹局域算符可以按某个重正化方案（例如
$\overline{MS}$ 方案）中的重正化算符作短距离展开。容易证明，该短
距离展开导致
\begin{equation}\label{eq:shortdist}
b_n^{(s,\mathrm{twist}-2)}\big(\sqrt{\tau}\Lambda_{\mathrm{QCD}}\big) =
\widetilde{C}_n^{(0)}(\sqrt{\tau}\mu)a^{(n)}(\mu) + {\cal
O}(\sqrt{\tau}\Lambda_{\mathrm{QCD}}),
\end{equation}
其中 $\mu$ 是重正化能标。该展开的
领头项 $a^{(n)}(\mu)$ 是重正化扭度2
算符的矩阵元，此算符与左端矩阵元中出现的涂抹
算符具有相同的 gamma 矩阵与导数结构。更高
阶的项来自进入涂抹矩阵元短距离展开的高维算符。

现在，把这个短距离展开与式~\eqref{eq:btwist} 结合起来写成
\begin{equation}\label{eq:double}
b_n^{(s)}\big(\sqrt{\tau}\Lambda_{\mathrm{QCD}}\big) =
C_n^{(0)}(\sqrt{\tau}\mu, \sqrt{\tau}P_z)a^{(n)}(\mu) + {\cal
O}\left(\sqrt{\tau}\Lambda_{\mathrm{QCD}},\frac{\Lambda_\mathrm{QCD}^2}{P_z^2}
\right),
\end{equation}
领头短距离系数函数
$ C_n^{(0)}(\sqrt{\tau}\mu, \sqrt{\tau}P_z)$ 以及各高阶修正
均可在微扰论中计算，因此该近似可以
系统地改进。
在本讨论的余下部分，我们假定工作在具有如下能标层级的区域：
\begin{equation}\label{eq:scales}
\Lambda_{\mathrm{QCD}},M_N\ll P_z\ll  \tau^{-1/2},
\end{equation}
从而幂修正与高扭度效应均可忽略。我们还
假定已施加靶质量修正。

为把涂抹准 PDF 与光前
PDF 联系起来，我们引入核函数 $Z(x,\sqrt{\tau}\mu,\sqrt{\tau}P_z)$，其 Mellin 矩由下式给出：
\begin{equation}\label{eq:cinvz}
 \left[C_n^{(0)}(\sqrt{\tau}\mu, \sqrt{\tau}P_z)\right]^{-1} =
\int_{-\infty}^{\infty} dx\, x^{n-1}
Z(x,\sqrt{\tau}\mu, \sqrt{\tau}P_z).
\end{equation}
利用这一定义以及乘法卷积的性质，我们发现
\begin{equation}
f(x,\mu) =  \int_{-\infty}^{\infty} \frac{d\xi}{\xi}\,
Z\left(\frac{x}{\xi},\sqrt{\tau}\mu, \sqrt{\tau}P_z\right)
q^{\,(s)}\left(\xi,\sqrt{\tau} \Lambda_{\mathrm{QCD}}\right) + {\cal
O}(\sqrt{\tau}\Lambda_{\mathrm{QCD}}).
\end{equation}
我们经由
\begin{equation}
C_n^{(0)}(\sqrt{\tau}\mu, \sqrt{\tau}P_z) = \int_{-\infty}^{\infty} dx\,
x^{n-1}
\widetilde{Z}(x,\sqrt{\tau}\mu, \sqrt{\tau}P_z),
\end{equation}
引入逆核，由此得到
\begin{equation}
q^{\,(s)}\left(x,\sqrt{\tau} \Lambda_{\mathrm{QCD}}, \sqrt{\tau}P_z\right) =
\int_{-1}^{1}
\frac{d\xi}{\xi}\, \widetilde{Z}\left(\frac{x}{\xi},\sqrt{\tau}\mu,
\sqrt{\tau}P_z\right)
f(\xi,\mu) +
{\cal O}(\sqrt{\tau}\Lambda_{\mathrm{QCD}})
\end{equation}
注意，以上所有关系仅在
\begin{equation}\label{eq:scales_tau}
\Lambda_{QCD},M_N\ll P_z \ll \tau^{-1/2}.
\end{equation}
时才有效。

核函数可以在连续时空微扰论中计算，所循方法由文献 \cite{Luscher:2011bx} 引入，例子见
文献~\cite{Ma:2014jla,Ji:2015jwa,Chen:2016fxx,Ishikawa:2016znu}。
鉴于那些计算是在连续时空中进行的，
它们与用于提取本文引入的
涂抹准 PDFs 的格点形式无关，而且涂抹准 PDFs 的
格点计算可以用多种格点作用量进行，其中每一种作用量都给出同一个普适的连续时空准 PDF。因此，格点
计算与把这些普适的
连续时空准 PDFs 匹配到光前 PDFs 的手续是彼此分离的。

在我们的方法中，一如在 Ji 的方法中，
大的核子动量只起压低高
扭度贡献的作用。然而，我们引入了一个新能标，即（逆）
流时间 $\tau^{-1}$，它充当紫外发散的调节子，这些发散来自定义矩阵元的复合外算符。能标
$\tau^{-1}$ 需要相对强子能标很大，但仍保持
有限。式~\eqref{eq:scales_tau} 所表达的
能标层级要求，与把物理截面因子化为 PDFs 与 Wilson 系数时的要求在本质上并无不同，
精神上类似于文献
\cite{Ma:2014jla,Ishikawa:2016znu} 提出的因子化方法。在该方法中，重正化
能标（扮演与我们流时间能标类似的角色）
和因子化能标相互独立，并与压低高扭度效应的大动量分开。
